\chapter{Introduction}
\label{ch:Introduction}

Commercial aviation has become integral to people’s lives over the past decade, primarily facilitating long-distance travel between airports and connecting cities worldwide. However, after landing at the airport, usually another form of transportation (taxi, car or public transport) is needed to travel to the final destination. This raises the question: what if the so-called \enquote*{last mile} can also be travelled by air?

This is where the concept of Inter-Urban Air Mobility (Inter-UAM) emerges. Inter-UAM envisions flying within large metropolitan areas or between smaller cities, offering a more extended range than Urban Air Mobility (UAM), greater than 100 km compared to approximately 30 km, and the capability to take off and land vertically at urban vertiports. 

This report aims to identify a promising Inter-UAM design, the \textit{Swing}, focused on minimising the aircraft’s ground and environmental footprint. The small ground footprint will enable the vehicle to take off from smaller vertiports, helipads or even driveways, allowing it to penetrate urban areas and land where other UAM vehicles cannot. Additionally, the vehicle will be fully electric, contributing to a greener future. The following Mission Need Statement (MNS) and Project Objective Statement (POS) define the project aims: 

\paragraph{MNS:} \textit{Achieve Sustainable Inter-Urban Air Mobility with a low ground footprint vehicle.}
\vspace{-5mm}
\paragraph{POS:} \textit{Design a low ground footprint, sustainable, urban transwing electric Vertical Take-Off and Landing vehicle within production costs of \qty[exponent-mode=input]{2}{M\texteuro}, by 10 students in 10 weeks.}

The present report can be split into three main parts: the first one focuses on analysing and refining the mission and on the preliminary design.
This is done through a market analysis in \Cref{ch:Market_Analysis} to assess the market, customer, and competition.
This is followed by an in-depth functional analysis for each mission phase and a risk analysis in  \Cref{ch:Functional_Analysis,ch:Technical_Risk_Assessment}.
The mission profile is defined and optimised in \Cref{ch:Mission_Profile} from the functions and market analysis.
The final chapter of the first part, \Cref{ch:Configuration}, presents the trade-off of different concepts and preliminary sizing, which introduces the concept that will be designed in detail.

The second part of this report focuses on the detailed conceptual design and analysis of the different subsystems of the \textit{Swing} aircraft.
In \Cref{ch:AerodynamicDesign,ch:Aerodynamic_Analysis}, the aerodynamic surfaces will be sized and designed, and the aerodynamic properties of the entire eVTOL will be analysed.
In \Cref{ch:Power}, the power curve for the vehicle is constructed, from which the propulsion can be sized in \Cref{ch:Propulsion}.
After this, the stability and control of the aircraft is analysed to design the empennage and control surfaces accordingly, in \Cref{ch:Stability_and_Control}.
Since the aircraft is an autonomous and electric VTOL, the autonomous control system and electric power system are also of great importance and are designed and analysed in \Cref{ch:Autonomous_System,ch:eps} respectively.
Finally, \Cref{ch:Structures_and_Materials} will analyse the aircraft’s structure and the hinge mechanism, where the materials will also be chosen.

The last part of this report will explain the resource allocation and interfaces of the design in \Cref{ch:Resource_Allocation_and_Budget_Breakdown,ch:Interfaces}.
In \Cref{ch:Operations_and_Logistics}, the operations and logistics of the aircraft will be described, followed by the sustainability approach in \Cref{ch:Sustainability_Approach}.
Afterwards, a financial analysis is conducted in \Cref{ch:Financial_Analysis}.
In \Cref{ch:Production_Plan}, the concept for the production of the aircraft will be sketched, followed by the future development of the project in \Cref{ch:Future_Development_Of_Project}.
As a conclusion to the report, in \Cref{ch:Requirements_Compliance} the final design will be checked for compliance with the original requirements, and an extensive compliance matrix is constructed accordingly.